\section{Algoritmus prekladu}
Po vygenerovaní stromu odvodenia prichádza na rad preklad do relačnej algebry. Preklad sa realizuje
v piatich krokoch. 

V prvom kroku využijeme dva objekty, ktoré nám automaticky vygeneroval ANTLR. Najprv vytvoríme
objekt \texttt{datalogLexer}, ktorý potom premeníme na \\\texttt{CommonTokenStream}. Tento druhý
objekt je v podstate zoznam tokenov, ktoré sa vyskytujú v zdrojovom kóde. Formálne by to boli znaky
(resp. písmená) slova, z ktorých hľadáme strom odvodenia.

\begin{lstlisting}
CharStream program;
datalogLexer lexer = new datalogLexer(program);
CommonTokenStream tokens = new CommonTokenStream(lexer);
\end{lstlisting}

V druhom kroku vytvoríme objekt \texttt{datalogParser} (ktorý tiež vygeneroval \\ANTLR), ktorý nám
umožní pracovať s gramatikou, ktorú sme definovali v súbore \texttt{datalog.g4}. Tento objekt nám
umožní vytvoriť strom odvodenia z tokenov, ktoré sme získali v prvom kroku.

\begin{lstlisting}
datalogParser parser = new datalogParser(tokens);
ParseTree tree = parser.program();
\end{lstlisting}

\subsection{Generácia AST}
Tretí krok je zavolanie objektu \texttt{DatalogASTBuilder}, ktorý nám transformuje strom odvodenia
na abstraktný syntaktický strom (AST). Tento strom je jednoduchší a pohodlnejší na ďalšie
spracovanie. Transformácia sa realizuje pomocou tzv. \textit{visitor patternu}. Vstupom je objekt
typu \texttt{DatalogProgram}, ktorý obsahuje všetky definície, klauzuly a dotaz. Triedy, ktoré sa
týkajú AST, sú definované v priečinkoch \texttt{datalog} a \texttt{common} (spoločné objekty pre
datalog aj relačnú algebru, termy a parametre).

\begin{lstlisting}
DatalogASTBuilder astBuilder = new DatalogASTBuilder();
DatalogProgram ast = (DatalogProgram) astBuilder.visit(tree);
\end{lstlisting}

Príkladom, aká transformácia sa deje, je transformácia na objekty typu \texttt{Atom} a \texttt{Predicate}.
Všetky nasledujúce podstromy odvodenia sa transformujú na rovnaké objekty s rovnakým rozhraním.

\begin{center}
    \[
    \Tree
    [.\texttt{predicate} 
        [.\texttt{normal\_predicate}
            \texttt{name}
            \texttt{(}
            \textit{possible terms}
            \texttt{)}
        ]
    ]
    \]
    \textit{Obyčajný predikát}
\end{center}

\begin{center}
    \[
    \Tree
    [.\texttt{predicate}
        [.\texttt{negative\_predicate}
            \qroof{\texttt{\textbackslash+}}.\texttt{NOT}
            [.\texttt{normal\_predicate}
                \texttt{name}
                \texttt{(}
                \textit{possible terms}
                \texttt{)}
            ]
        ]
    ]
    \]
    \textit{Negovaný obyčajný predikát}
\end{center}

\begin{center}
    \[
    \Tree
    [.\texttt{built\_in\_predicate} 
        [.\texttt{normal\_built\_in\_predicate}
            \qroof{\texttt{\$}}.\texttt{BUILT\_IN}
            \texttt{name}
            \texttt{(}
            \textit{possible parameters}
            \texttt{)}
        ]
    ]
    \]
    \textit{Vstavaný predikát}
\end{center}

\begin{center}
    \[
    \Tree
    [.\texttt{built\_in\_predicate}
        [.\texttt{negative\_built\_in\_predicate}
            \qroof{\texttt{\textbackslash+}}.\texttt{NOT}
            [.\texttt{normal\_built\_in\_predicate}
                \qroof{\texttt{\$}}.\texttt{BUILT\_IN}
                \texttt{name}
                \texttt{(}
                \textit{possible parameters}
                \texttt{)}
            ]
        ]
    ]
    \]
    \textit{Negovaný vstavaný predikát}
\end{center}

Celá štruktúra a hierarchia tried, ktoré reprezentujú AST, sú podrobnejšie popísané v nasledujúcom
diagrame.

\begin{figure}[h]
\centering
\resizebox{\linewidth}{!}{%
\begin{forest}
  for tree={
    draw,
    align=left,
    font=\ttfamily\footnotesize,
    inner sep=5pt,
    fill=gray!10,
    edge={-{Stealth[scale=0.8]}},
    l sep=1.2cm,
    s sep=0.3cm,
  }
  [{\bfseries DatalogProgram\\ - clauses: List<Clause>\\ - definitions: List<Definition>\\ - query: Query}
    [{\bfseries Clause\\ - head: Predicate\\ - body: List<Atom>}
      [{\bfseries Predicate\\ - name: String\\ - terms: List<Term>\\ - isBuiltIn: bool}
        [{\itshape\bfseries Term}
          [{\bfseries Constant\\ - value: String},       edge={dashed,-{Triangle[open,scale=0.8]}}]
          [{\bfseries Variable\\ - name: String},        edge={dashed,-{Triangle[open,scale=0.8]}}]
          [{\bfseries FunctionTerm\\ - name: String\\ - args: List<Term>}, edge={dashed,-{Triangle[open,scale=0.8]}}]
          [{\bfseries Parameter\\ - text: String},       edge={dashed,-{Triangle[open,scale=0.8]}}]
        ]
      ]
      [{\bfseries Atom\\ - predicate: Predicate\\ - isNegated: boolean}]
    ]
    [{\bfseries Definition\\ - name: String\\ - arity: int\\ - tuples: List<List<Term>>}]
    [{\bfseries Query\\ - predicate: Predicate}]
  ]
\end{forest}%
}
\caption*{Hierarchia tried AST}
\end{figure}

V tejto fáze sa tiež realizuje kontrola, či sú klauzuly bezpečné a či všetky \\\verb|term_tuple| v
definícii majú rovnakú aritu, keďže tieto vlastnosti programu nie je možné vyjadriť pomocou
bezkontextovej gramatiky. Taktiež sa kontroluje, či termy v \verb|term_tuple| neobsahujú premenné
(táto vlastnosť by sa síce dala vyjadriť bezkontextovou gramatikou, ale uvedený postup je
jednoduchší).

\subsection{Preklad do relačnej algebry}
Posledné 2 kroky sa týkajú samotného prekladu do relačnej algebry a vykonáva ich trieda
\texttt{DatalogCompiler}.

\begin{lstlisting}
DatalogCompiler compiler = new DatalogCompiler();
RAExpr compiled = compiler.compile(ast);
\end{lstlisting}

V štvrtom kroku vygenerujeme pre každý predikát v programe objekt typu \texttt{RAExpr}, ktorý
reprezentuje reláciu k nemu prislúchajúcu. Ak je nejaký predikát závislý na inom predikáte, zatiaľ
ho nepreložíme, ale použijeme referenciu (\textit{placeholder}) na tento predikát (objekt typu
\texttt{RelationRef}). Prekladať nebudeme priamo do relačnej algebry, ale do Java objektov
reprezentujúcich relačnú algebru. To nám umožní ľahšie manipulovať s výsledným výrazom a taktiež
jednoduchšie meniť zápis výsledného výrazu. Napríklad ak by sme chceli zmeniť poradie parametrov v
operátoroch, alebo zmeniť odsadenie a zátvorky.

\vline

Prvým problémom, s ktorým sa v tomto kroku stretávame, je ako preložiť definície. Definície sú v
podstate relácie, ktoré obsahujú pevne dané množiny n-tíc. Keďže nemáme žiaden operátor, ktorý by
toto priamo podporoval, musíme \uv{zneužiť} operátor \JOIN\ a vstavané relácie.

\begin{definition}
    Vstavané relácie $\$true$ a $\$false$ sú definované nasledovne:
    \begin{align*}
        \$true &= \{()\} \\
        \$false &= \emptyset
    \end{align*}
\end{definition}

Zoberme si výraz $\JOIN([\$true], [[]], [n1(), a()])$ (normálne by sme zapísali konštanty ako $1$ a
$a$, ale naša aktuálna implementácia vyžaduje tento zápis). Nech $M$ je množina z definície \JOIN.
Potom $M = \{()\}$, keďže iba pre $()$ platí $() \in \$true$. Zápis je trochu neintuitívny, ale robí
čo potrebujeme. Výsledná relácia $R$ je rovná $\{(n1(), a())\}$.

Týmto spôsobom vieme vyrobiť ľubovoľnú reláciu s jednou
n-ticou. Pre viacnásobné n-tice môžeme použiť viacnásobné \JOIN\ \textit{obalené} v \UNION.

\begin{priklad}
    Program (resp. iba časť programu) \verb|color :- {(red()), (blue())}.| sa preloží ako
\begin{verbatim}
UNION([
    JOIN([$true], [[]], [red()]),
    JOIN([$true], [[]], [blue()])
])
\end{verbatim}
\end{priklad}

Keďže jeden predikát môže byť definovaný viackrát, musíme prejsť všetky definície, pozbierať všetky
n-tice a spojiť ich do jednej relácie (bez duplikátov). 

Pôvodne sme chceli mať v relačnej algebre aj operátor \DEF, ktorý by nám umožňoval definovať reláciu
pomocou množiny n-tíc. Preto sa používa v programe kompilátora trieda \texttt{RADef}. Keďže sme sa
ale nakoniec rozhodli uprednostniť menšiu množinu operátorov, \DEF\ sa transformuje na \UNION\
\JOIN-ov, ako sme ukázali v príklade vyššie. Preto sa v tejto fáze realizuje aj táto transformácia.

\vline

Druhým problémom je preklad klauzúl. V Datalogu je klauzula v podstate konjunkcia atómov a keďže
konjunkcia je asociatívna a komutatívna, môžeme atómy v klauzule zoradiť ľubovoľne. Preto môžeme
pred prekladom klauzúl zoradiť atómy tak, že najprv budú kladné atómy a až potom negované atómy.

Z kladných atómov vieme pomocou \JOIN\ vyrobiť reláciu, ktorá obsahuje všetky n-tice, ktoré môžeme
doplniť, aby boli všetky pravdivé. Túto výslednú reláciu potom použijeme ako prvú
reláciu v \ANTIJOIN\ operátore, pomocou ktorého odfiltrujeme všetky n-tice, ktoré by spĺňali
negované atómy. Výstupom celej klauzuly sú ale iba tie n-tice, ktoré spĺňajú hlavu klauzuly. To
zohľadníme vo výstupnom vektore v \ANTIJOIN-e (respektíve v \JOIN-e, ak nemáme negatívne atómy).

\vline

\refstepcounter{priklad}
Časť Datalogového programu
\begin{verbatim}
dvojica(f(A), D) :- $hrana(A, B),
                    $hrana(f(B), kon()), 
                    $hrana(kon(), D),
                    \+ $hrana(A, kon()),
                    \+ $hrana(f(B), D).
\end{verbatim}

sa preloží nasledovne. Najprv vytvoríme reláciu z kladných atómov pomocou \JOIN-u.

\begin{verbatim}
JOIN(
    [$hrana, $hrana, $hrana],
    [[A, B], [f(B), kon()], [kon(), D]],
    [A, B, D]
)
\end{verbatim}

Táto relácia obsahuje všetky n-tice, ktoré spĺňajú kladné atómy. Potom použijeme \ANTIJOIN\ na
vytvorenie výslednej relácie \verb|dvojica|.

\begin{verbatim}
ANTIJOIN(
    [JOIN(
        [$hrana, $hrana, $hrana],
        [[A, B], [f(B), kon()], [kon(), D]],
        [A, B, D]
    ),
    $hrana, $hrana],
    [[A, B, D], [A, kon()], [f(B), D]],
    [f(A), D]
)
\end{verbatim}
\begin{center}
\centering
\textit{Príklad \thepriklad}
\end{center}
\medskip

V prípade, že klauzula neobsahuje žiadne negované atómy, použijeme iba \JOIN\ a výstupný vektor bude
zodpovedať hlave klauzuly. V prípade, že klauzula obsahuje iba jeden pozitívny atóm, \JOIN\ nie je
potrebné vyrábať, ale môžeme použiť tento atóm priamo v \ANTIJOIN-e.

Zaujímavý prípad je klauzula, ktorá obsahuje iba negované atómy. Keďže musí byť bezpečná, nemôže
obsahovať žiadne premenné, ale iba konštanty. V takomto prípade má klauzula tvar \uv{\textit{n-tica 1
patrí do relácie A, ak n-tica 2 nepatrí do relácie B, n-tica 3 nepatrí do relácie C, ...}}. Keďže
ale \ANTIJOIN\ potrebuje jednu reláciu v pozitívnom kontexte, musíme do všetkých takýchto klauzúl
pridať nejaký pozitívny atóm, napríklad \verb|$true()|.

\begin{priklad}
    Klauzula \verb|test(n1()) :- \+ $a(n2()), \+ $b(n3()).| sa preloží ako:
\begin{verbatim}
ANTIJOIN([$true, $a, $b], [[], [n2()], [n3()]], [n1()])
\end{verbatim}
\end{priklad}

V prípade, že jeden predikát je definovaný viackrát, musíme všetky klauzuly, ktoré mu prislúchajú,
spojiť do jednej relácie pomocou \UNION-u.

\begin{priklad}
    Program
\begin{verbatim}
p(X) :- $a(X).
p(Y) :- $b(Y).
\end{verbatim}
sa preloží ako:
\begin{verbatim}
UNION([
    JOIN([$a], [[X]], [X]),
    JOIN([$b], [[Y]], [Y])
])
\end{verbatim}
\end{priklad}

Výsledkom štvrtého kroku je teda mapa \texttt{Map<String, RAExpr>}, ktorá priraďuje každému
predikátu (resp. relácii) výraz v relačnej algebre, ktorý ho reprezentuje. Ak sa v niektorej relácii
vyskytuje referencia na iný predikát, táto referencia je zatiaľ nevyplnená (je to objekt typu
\texttt{RelationRef}).

\begin{note}
    Teoreticky nezáleží od poradia atómov v klauzule, ale pri konkrétnych vstavaných predikátoch
    môže záležať. Zoberme si napríklad predikát \\\verb|$plus(X, Y, Z)|, ktorý je pravdivý vtedy a
    len vtedy, keď $X + Y = Z$. Jeho implementácia v databázovom systéme môže byť taká, že ak sú
    zadané (sú konštanty) $X$ a $Y$, vráti $Z$ ($Z$ je premenná) alebo ak sú zadané $X$, $Y$ a $Z$,
    vráti pravdivostnú hodnotu a to je všetko. Aktuálna implementácia napr. \JOIN-u v našom
    databázovom systéme je ale nasledovná. Zoberie sa prvá relácia a jej vstupný vektor. Vyberie sa
    z nej nejaká n-tica, aby spĺňala vstupný vektor. Týmto sa môžu konkretizovať niektoré premenné.
    Potom sa zoberie druhá relácia a jej vstupný vektor (s už konkretizovanými premennými) a spraví
    sa to isté. Preto v našom prípade záleží od poradia atómov. Ako príklad si zoberme program
\begin{verbatim}
    p(Z) :- $cisloDo10(X), $cisloDo10(Y), $plus(X, Y, Z).
\end{verbatim}
    Príslušný výraz v relačnej algebre je
\begin{verbatim}
JOIN(
    [$cisloDo10, $cisloDo10, $plus],
    [[X], [Y], [X, Y, Z]],
    [Z]
)
\end{verbatim}
    Sémanticky identický program
\begin{verbatim}
p(Z) :- $plus(X, Y, Z), $cisloDo10(X), $cisloDo10(Y).
\end{verbatim}
    sa ale preloží ako
\begin{verbatim}
JOIN(
    [$plus, $cisloDo10, $cisloDo10],
    [[X, Y, Z], [X], [Y]],
    [Z]
)
\end{verbatim}
    a tento výraz nemusí byť správne vyhodnotený, keďže $plus$ nemusí vedieť pracovať s troma
    premennými.
    
    Preto si používateľ musí dávať pozor na to, v akom poradí píše atómy do klauzúl, keďže to môže
    ovplyvniť vypočítateľnosť dotazu. Náš kompilátor zoberie najprv kladné atómy v poradí a následne
    negované atómy v poradí v akom boli zapísané v programe.
\end{note}

\vline

V piatom kroku tieto referencie nahradíme výrazom, ktorý im prislúcha. Keďže
ale môže nastať situácia, že niektorý predikát je závislý sám na sebe (priamo alebo nepriamo),
musíme prípadne predikáty \textit{zabaliť} do rekurzívneho operátora \REC. To spravíme tak, že najprv
skonštruujeme graf závislosti predikátov, v ktorom bude orientovaná hrana z predikátu $A$ do
predikátu $B$, ak $A$ (priamo) závisí na $B$. Potom pre každý predikát skontrolujeme, či existuje
cesta z neho do seba samého. Ak áno, tento predikát je rekurzívny a musíme ho zabaliť do \REC.

Výsledný výraz v relačnej algebre začneme stavať od dotazu. Prechádzame jeho výraz do hĺbky a
pamätáme si, cez aké \REC\ operátory sme prešli. Keď narazíme na \texttt{RelationRef} (pre ktorý sme
nevideli operátor \REC), nahradíme ho výrazom, ktorý mu prislúcha. Takto postupujeme, až kým
nenahradíme všetky \texttt{RelationRef} výrazmi, ktoré im prislúchajú. Výsledkom ale nemá byť celá
relácia, ale iba tie záznamy, ktoré spĺňajú hlavu dotazu. Preto ešte celý výsledný výraz zabalíme do
\JOIN-u, ktorého výstupný vektor bude zodpovedať hlave dotazu.

\medskip

\refstepcounter{priklad}
    Ukážeme si preklad programu spolu s vysvetlením, ktoré časti relačnej algebry zodpovedajú ktorým
    častiam programu.
\begin{verbatim}
 1  node :- {(a()), (b()), (c())}.
 2  blocked :- {(c())}.
 3  
 4  reachable(X) :- node(X).
 5  reachable(X) :- reachable(X).
 6  
 7  selected(X) :- reachable(X), $eq(X, b()).
 8  free(X) :- reachable(X), \+ blocked(X).
 9  
10  result(X) :- free(X), selected(X).
11  
12  ?- result(Z).
\end{verbatim}

Znak \texttt{\%} označuje jednoriadkový komentár.

\begin{verbatim}
JOIN([                                    % dotaz - 12
  JOIN([                                  % result - 10
    ANTIJOIN([                            % free - 8
      REC(reachable,                      % lebo riadok 5 je rekurzívny
        UNION([
          JOIN([                          % reachable - 4
            UNION([                       % node - 1
              JOIN([$true], [[]], [a()]),
              JOIN([$true], [[]], [b()]),
              JOIN([$true], [[]], [c()])
            ])
          ], [[X]], [X]),
          JOIN([reachable], [[X]], [X])   % reachable - 5
        ])
      ),
      JOIN([$true], [[]], [c()])          % blocked - 2
    ], [[X], [X]], [X]),
    JOIN([                                % selected - 7
      REC(reachable,                      % znovu reachable
        UNION([
          JOIN([
            UNION([
              JOIN([$true], [[]], [a()]),
              JOIN([$true], [[]], [b()]),
              JOIN([$true], [[]], [c()])
            ])
          ], [[X]], [X]),
          JOIN([reachable], [[X]], [X])])
      ),
      $eq
    ], [[X], [X, b()]], [X])
  ], [[X], [X]], [X])
], [[Z]], [Z])
\end{verbatim}
\begin{center}
\centering
\textit{Príklad \thepriklad}
\end{center}
\medskip

\subsection*{Rekapitulácia}
Preklad do relačnej algebry sa realizuje v piatich krokoch. Heslovite sú to:
\begin{enumerate}
    \item Vytvorenie lexeru a token streamu
    \item Vytvorenie parseru a stromu odvodenia
    \item Generácia AST z odvodenia
    \item Preklad predikátov do relačnej algebry s nevyplnenými referenciami
    \begin{enumerate}
        \item Preklad definícií pomocou \JOIN\ a \UNION\
        \item Preklad klauzúl pomocou \JOIN\ a \ANTIJOIN\
        \item Spojenie klauzúl do jednej relácie pomocou \UNION\
    \end{enumerate}
    \item Zostavenie výsledného výrazu v relačnej algebre
    \begin{enumerate}
        \item Zostavenie grafu závislosti predikátov
        \item Zabalenie rekurzívnych predikátov do \REC\
        \item Nahradenie referencií výrazmi, ktoré im prislúchajú
        \item Zabalenie výsledného výrazu do \JOIN-u reprezentujúceho dotaz
    \end{enumerate}
\end{enumerate}